% part: intuitionistic-logic
% chapter: propositions-as-types
% section: normalization

\documentclass[../../../include/open-logic-section]{subfiles}

\begin{document}

\olfileid[pt]{int}{pty}{nor}

\olsection{Normalização}

Nesta seção, demonstramos que, segundo certa ordem de redução, qualquer
dedução pode ser reduzida a uma dedução normal; isso é chamado de
\emph{propriedade de normalização}. Usaremos a correspondência entre
proposições e tipos: mostramos que todo termo de prova pode ser reduzido
a uma forma normal; disso segue a normalização para !!{derivation}s de
dedução natural.

Primeiro, definimos algumas funções que medem a complexidade dos termos.
O \emph{comprimento}~$\len{!A}$ de !!a{formula}s é definido por
\begin{align*}
  \len{p} &= 0\\
  \len{!A \land !B} &= \len{!A} + \len{!B} + 1\\
  \len{!A \lor !B} &= \len{!A} + \len{!B} + 1 \\
  \len{!A \lif !B} &= \len{!A} + \len{!B} + 1.
\end{align*}
A complexidade de um redex~$M$ é medida por seu \emph{posto de
  corte}~$\cutrank{M}$:
\begin{align*}
  \cutrank{(\lambd[\typeof{x}{!A}][\typeof{N}{!B}])Q} &=
  \len{!A} + \len{!B} + 1 \\
  \cutrank{\ande{i}{\andi{\typeof{M}{!A}}{\typeof{N}{!B}}}} &=
  \len{!A} + \len{!B} + 1 \\
  \cutrank{\ore{\ori{i}{!A_{i}}{\typeof{M}{!A_i}}}
    {\typeof{x_1}{!A_1}}{\typeof{N_1}{!C}}
    {\typeof{x_2}{!A_2}}{\typeof{N_2}{!C}}} &=
  \len{!A} + \len{!B} + 1
\end{align*}
A complexidade de um termo de prova é medida pelo redex mais complexo
que ele contém, e é $0$ se ele é normal:
\[
  \maxrank{M} = \max \{\cutrank{N} | N \text{ é um subtermo de } M \text{
    e é um redex}\}
\]

\begin{lem}\ollabel{lem:subst}
  Se $\Subst{M}{\typeof{N}{!A}}{\typeof{x}{!A}}$ é um redex e $M
  \not\ident x$, então vale um dos seguintes casos:
  \begin{enumerate}
  \item o próprio $M$ é um redex; ou
  \item $M$ tem a forma $\ande{i}{x}$, e $N$ tem a forma
    $\andi{P_1}{P_2}$
  \item $M$ tem a forma $\ore{i}{x_1}{P_1}{x_2}{P_2}$, e $N$ tem
    a forma $\ori{i}{}{Q}$
  \item $M$ tem a forma $x Q$, e $N$ tem a forma
    $\lambd[x][P]$
  \end{enumerate}
  No primeiro caso, $\cutrank{\Subst{M}{N}{x}} = \cutrank{M}$; nos
  outros casos, $\cutrank{\Subst{M}{N}{x}} = \len{!A})$.
\end{lem}
\begin{proof}
  Demonstração por indução em~$M$.
  \begin{enumerate}
  \item Se $M$ é uma única variável $y$ e $y \not\ident x$, então
    $\Subst{y}{N}{x}$ é $y$ e, portanto, não é um redex.
  \item Se $M$ tem a forma $\andi{N_1}{N_2}$, $\lambd[x][N]$ ou
    $\ori{i}{!A}{N}$, então $\Subst{M}{\typeof{N}{!A}}{\typeof{x}{!A}}$
    também tem essa forma e, portanto, não é um redex.
  \item Se $M$ tem a forma $\ande{i}{P}$, consideramos dois casos.
    \begin{enumerate}
      \item Se $P$ tem a forma $\andi{P_1}{P_2}$, então $M \ident
        \ande{i}{\andi{P_1}{P_2}}$ é um redex e, claramente,
        \[
        \Subst{M}{N}{x} \ident
        \ande{i}{\andi{\Subst{P_1}{N}{x}}{\Subst{P_2}{N}{x}}}
        \]
        também é um redex. Os postos de corte são iguais.
      \item Se $P$ é uma única variável, ela deve ser $x$ para que a
        substituição seja um redex, e $N$ deve ter a forma
        $\andi{P_1}{P_2}$. Considere agora
        \[
        \Subst{M}{N}{x} \ident \Subst{\ande{i}{x}}{\andi{P_1}{P_2}}{x},
        \]
        que é $\ande{i}{\andi{P_1}{P_2}}$. Seu posto de corte é igual a
        $\cutrank{x}$, que é $\len{!A}$.
    \end{enumerate}
  \end{enumerate}
  Os casos de $\ore{N}{x_1}{N_1}{x_2}{N_2}$ e $P Q$ são semelhantes.
\end{proof}

\begin{lem}
  Se $M$ se contrai a $M'$ e $\cutrank{M} > \cutrank{N}$ para todos os
  subtermos redex próprios $N$ de $M$, então $\cutrank{M} >
  \maxrank{M'}$.
\end{lem}

\begin{proof}
  Demonstração por casos.
  \begin{enumerate}
  \item Se $M$ tem a forma $\ande{i}{\andi{M_1}{M_2}}$, então $M'$ é
    $M_i$; como todo subtermo de $M_i$ também é um subtermo próprio de
    $M$, vale a afirmação.
  \item Se $M$ tem a forma
    $(\lambd[\typeof{x}{!A}][N])\typeof{Q}{!A}$, então $M'$ é
    $\Subst{N}{\typeof{Q}{!A}}{\typeof{x}{!A}}$. Considere um redex em
    $M'$. Ou há um redex correspondente em $N$ com o mesmo posto de
    corte, que é menor do que $\cutrank{M}$ por hipótese, ou o posto de
    corte é igual a $\len{!A}$, que, por definição, é menor do que
    $\cutrank{(\lambd[x^{!A}][N])Q}$.
  \item Se $M$ tem a forma
    \[
    \ore{\ori{i}{}{\typeof{N}{!A_i}}}
        {\typeof{x_1}{!A_1}}{\typeof{N_1}{!C}}
        {\typeof{x_2}{!A_2}}{\typeof{N_2}{!C}},
    \]
    então $M' \ident \Subst{N_i}{N}{\typeof{x_i}{!A_i}}$. Considere um
    redex em~$M'$. Ou há um redex correspondente em $N_i$ com o mesmo
    posto de corte, que é menor do que $\cutrank{M}$ por hipótese, ou o
    posto de corte é igual a $\len{!A_i}$, que, por definição, é menor do que
    $\cutrank{\ore{\ori{i}{}{\typeof{N}{!A_i}}}
      {\typeof{x_1}{!A_1}}{\typeof{N_1}{!C}}
      {\typeof{x_2}{!A_2}}{\typeof{N_2}{!C}}}$.
  \end{enumerate}
\end{proof}

\begin{thm}
  Todos os termos de prova se reduzem à forma normal; todas as
  !!{derivation}s se reduzem a !!{derivation}s normais.
\end{thm}

\begin{proof}
  A segunda afirmação segue da primeira. Demonstramos a primeira por
  indução completa em $m=\maxrank{M}$, em que $M$ é um termo de prova.
  \begin{enumerate}

  \item Se $m=0$, $M$ já é normal.
  \item
    Caso contrário, procedemos por indução em~$n$, o número de
    redexes em $M$ cujo posto de corte é igual a~$m$.
    \begin{enumerate}
    \item Se $n=1$, selecione qualquer redex $N$ tal que $m = \cutrank{N} >
      \cutrank{P}$ para todo subtermo próprio $P$ que também seja, é
      claro, um redex. Tal redex deve existir, pois todo termo tem
      apenas um número finito de subtermos.

      Seja $N'$ o reduto de $N$. Pelo lema,
      $\maxrank{N'} < \maxrank{N}$; vemos, portanto, que $n$, o número
      de redexes com $\cutrank=m$, diminui. Assim, $m$ diminui (em $1$
      ou mais), e podemos aplicar a hipótese de indução para~$m$.
    \item Para o passo de indução, suponha $n>1$. O processo é
      semelhante, exceto que $n$ diminui apenas até um número positivo
      e, portanto, $m$ não muda. Aplicamos simplesmente a hipótese de
      indução para~$n$.
    \end{enumerate}
\end{enumerate}
\end{proof}

O fato de os $\lambd$-termos tipados se normalizarem significa que todo
termo desse tipo \emph{tem uma forma normal}. Se entendermos a redução
$\beta$ de $\lambd$-termos à sua forma normal como a execução de uma
computação e a forma normal como o valor da computação, isso significa
que toda computação no $\lambd$-cálculo tipado acaba por parar e produzir
um valor. Além disso, a preservação de tipos garante que o valor tenha
o tipo correto. Por exemplo, se $N$ tem tipo $!A \lif !B$ (isto é, é uma
função que deve operar sobre argumentos do tipo~$!A$ e produzir valores
do tipo~$!B$) e o aplicamos a um termo~$M$ do tipo~$!A$ (a entrada),
$(NM)$ se reduz a um termo $N'$ (a saída), e essa saída é garantidamente
do tipo~$!B$.

Na verdade, os termos do $\lambd$-cálculo tipado não apenas se reduzem a
formas normais segundo uma maneira de aplicar reduções~$\beta$ a
subtermos: \emph{qualquer} maneira de aplicar reduções~$\beta$ a
subtermos acaba por terminar em uma forma normal. Essa propriedade é
chamada de \emph{normalização forte}. Ela significa não apenas que há
uma maneira de garantir que uma computação acabe por parar e produzir um
valor, mas que \emph{qualquer} maneira de executar a computação acaba por
produzir um valor.

Como sabemos que diferentes maneiras de executar uma computação produzem
o \emph{mesmo} valor? Até agora, sabemos apenas (pela preservação de
tipos) que todos os valores que ela produz têm o mesmo tipo. O
$\lambd$-cálculo tipado tem outra propriedade importante: ele é
\emph{confluente}, ou \emph{Church--Rosser}. Se $N \red N_1$ e $N \red
N_2$, então existe algum termo~$N'$ tal que $N_1 \red N'$ e $N_2 \red
N'$. Lembre que $N \red N'$ significa que $N'$ resulta de zero ou mais
passos de redução~$\redone$. Se $N_1$ está na forma normal, então o único
termo $N'$ tal que $N_1 \red N'$ é o próprio $N_1$ (por meio de zero
passos de redução~$\redone$). Segue-se que, se $N_1$ e $N_2$ estão ambos
na forma normal, então $N_1 = N' = N_2$. Em outras palavras, como
$\red$ é fortemente normalizante, toda maneira de reduzir um termo acaba
por resultar em uma forma normal. Como $\red$ é Church--Rosser, quaisquer
duas formas normais são iguais. Portanto, a computação no
$\lambd$-cálculo tipado sempre termina, e o valor (a forma normal) é o
mesmo, independentemente de como a computação procede.

Demonstrar que uma relação é confluente costuma ser difícil. No entanto,
podemos estabelecer facilmente a seguinte condição mais fraca:

\begin{prop}[Propriedade fraca de Church--Rosser]
Se $N \redone N_1$ e $N \redone N_2$, então existe um termo $N'$ tal
que $N_1 \red N'$ e $N_2 \red N'$.
\end{prop}

\begin{proof}
  Os dois termos $N_1$ e $N_2$ resultam de $N$ pela substituição,
  respectivamente, do redex $M_1$ ou $M_2$ por seu reduto. Os redexes
  $M_1$ e $M_2$ são subtermos de $N$ e, portanto, não podem se
  sobrepor. Isso significa que ou um é subtermo do outro, ou eles
  ocorrem em posições diferentes e não sobrepostas de $N$. Considere o
  segundo caso: $N$ tem a forma $N(M_1,M_2)$. Substituir $M_1$ por seu
  reduto $M_1'$ resulta em $N_1$, que é $N(M_1',M_2)$, e substituir
  $M_2$ por seu reduto $M_2'$ resulta em $N_2$, que é $N(M_1,M_2')$.
  Ambos se reduzem a $N(M_1',M_2')$. No outro caso, suponha que $M_2$
  seja um subtermo de~$M_1$ (o outro caso é simétrico) e distinga casos
  segundo o tipo de redex que $M_1$ é. Por exemplo, se $M_1 \ident
  \proj{1}{\pair{O_1}{O_2}}$, ele se reduz a $O_1$. Se $M_2$ é um
  subtermo de $O_1$, então converter $M_2$ resulta em $O_1'$. Assim,
  converter primeiro $M_1$ e depois $M_2$ resulta em $O_1'$. Mas
  converter primeiro $M_2$ resulta em $\proj{1}{\pair{O_1'}{O_2}}$, e
  isso também se reduz a~$O_2$.
\end{proof}

\begin{lem}[Lema de Newman]
Se $\red$ é fortemente normalizante e tem a propriedade fraca de
Church--Rosser, então é confluente.
\end{lem}

\begin{proof}
  Como $\red$ é fortemente normalizante, toda sequência de redução
  resulta em uma forma normal. As formas normais são únicas sse $\red$
  é confluente. Suponha, então, que algum termo $N$ tenha duas formas
  normais distintas $N'$ e $N''$, resultantes das sequências de redução
  \begin{align*}
  N & = N_1' \redone N_2' \redone \dots \redone N_k' = N' \text{ e}\\
  N & = N_1'' \redone N_2'' \redone \dots \redone N_l'' = N''
  \end{align*}
  (A sequência de redução deve ter comprimento $>0$, pois, caso
  contrário, $N$ já estaria na forma normal.)

  Se $N_1' = N_1''$, então ele tem duas formas normais distintas ($N'$ e
  $N''$). Se $N_1' \neq N_1''$, então, pela propriedade fraca de
  Church--Rosser, existe algum $N'''$ tal que $N_1' \red N'''$ e
  $N_1'' \red N'''$. Como $N' \neq N''$, ou $N''' \neq N'$ ou $N'''
  \neq N''$. Consequentemente, ou $N_1'$ ou $N_1''$ tem duas formas
  normais distintas. Mostramos que, se $N$ tem duas formas normais
  distintas, existe um $N_1$ tal que $N \redone N_1$ e $N_1$ tem duas
  formas normais distintas. Evidentemente, podemos continuar e produzir
  $N \redone N_1 \redone N_2 \redone \dots$, mas isso é impossível,
  pois $\redone$ é fortemente normalizante.
\end{proof}


\end{document}
